\worksheettitle
  {Домашняя работа}
  {\modulemark{M02-H} Ноль и разрядный состав}

\studentnameblock

\begin{notebox}[Маршрут]
  Сначала определяйте разряды справа налево. При разложении не записывайте
  нулевые слагаемые. При сборке сохраняйте каждый разряд с помощью цифры \(0\).
  Если это вызывает затруднение, воспользуйтесь карточкой M02-R.
\end{notebox}

\begin{practicebox}[1. Заполните пропуски]
  \begin{enumerate}[label=\textbf{\arabic*.}]
    \item \(72\,305=70\,000+\answerblank[18mm]+300+\answerblank[16mm].\)
    \item \(90\,408=90\,000+\answerblank[18mm]+8.\)
    \item \(603\,070=600\,000+\answerblank[18mm]+70.\)
  \end{enumerate}
\end{practicebox}

\begin{practicebox}[2. Разложите на разрядные слагаемые]
  \begin{enumerate}[label=\textbf{\arabic*.}]
    \item \(48\,206=\answerblank[88mm].\)
    \item \(70\,090=\answerblank[88mm].\)
    \item \(305\,004=\answerblank[88mm].\)
    \item \(800\,060=\answerblank[88mm].\)
  \end{enumerate}
\end{practicebox}

\begin{practicebox}[3. Соберите числа]
  \begin{enumerate}[label=\textbf{\arabic*.}]
    \item \(40\,000+5\,000+300+2=\answerblank[30mm].\)
    \item \(90\,000+600+7=\answerblank[30mm].\)
    \item \(500\,000+20\,000+40+8=\answerblank[30mm].\)
    \item \(700\,000+3=\answerblank[30mm].\)
  \end{enumerate}
\end{practicebox}

\begin{warningbox}[4. Найдите и исправьте ошибки]
  \begin{enumerate}[label=\textbf{\arabic*.}]
    \item \(60\,405=60\,000+4\,000+5.\)
    \item \(80\,000+900+6=896.\)
    \item \(402\,070=400\,000+20\,000+70.\)
  \end{enumerate}
  Для каждой ошибки запишите причину и верный ответ.
  \writinglines{3}
\end{warningbox}

\begin{practicebox}[5. Сравните записи]
  \begin{enumerate}[label=\textbf{\arabic*.}]
    \item Чем отличаются числа \(50\,304\) и \(5\,304\)?
    \item Как изменилось значение цифры \(5\)?
    \item Почему число \(50\,304\) нельзя записать как \(5\,304\)?
  \end{enumerate}
  \writinglines{2}
\end{practicebox}

\begin{practicebox}[6. Составьте число]
  Запишите шестизначное число, в котором:
  \begin{itemize}
    \item цифра \(6\) имеет значение \(600\,000\);
    \item цифра \(4\) имеет значение \(4\,000\);
    \item цифра \(8\) имеет значение \(80\);
    \item остальные разряды содержат нули.
  \end{itemize}
  Число: \answerblank[35mm].

  Разложение: \answerblank[75mm].
\end{practicebox}

\begin{checkpointbox}[7. Объясните правило]
  Закончите предложение: «При записи числа нельзя пропускать разряд, даже если в нем нет единиц (или отсутствует слагаемое), потому что…»
  \writinglines{1}
  \textbf{Перед сдачей:} проверьте значения цифр в разложениях и убедитесь, что ни один разряд не пропущен при записи чисел.
\end{checkpointbox}
